Sequences in this chapter are indexed by $\N=\{0,1,2,\ldots\}$.
The range $\{u_n:n\in\N\}$ records values; ``finitely many exceptional terms'' refers to indices, even when values repeat.

\begin{definition} \label{an1:def:sequence-convergence}
  A \textbf{sequence} in $X$ is a function $u:\N\to X$, with terms $u_n=u(n)$.
  It \textbf{converges to} $p\in X$ if for every $\varepsilon>0$ there is $N\in\N$ such that $d(u_n,p)<\varepsilon$ whenever $n\geq N$.
  We write $u_n\to p$ or $\lim_{n\to\infty}u_n=p$.
  A sequence is \textbf{convergent} if it has a limit in its ambient space and \textbf{divergent} otherwise.
\end{definition}

% Lean source: Textbooks.MathematicalAnalysis1.Chapter02.ConvergentSequences.tendsto_iff
\begin{lemma} \label{an1:lem:tendsto-iff}
  The standard sequence limit is equivalent to the epsilon condition above.

  \begin{lean}
    theorem tendsto_iff (u : ℕ → X) (p : X) :
        Filter.Tendsto u Filter.atTop (nhds p) ↔
          ∀ ε : ℝ, 0 < ε → ∃ N : ℕ, ∀ n ≥ N, dist (u n) p < ε
  \end{lean}
\end{lemma}

\begin{proof}
  Balls form a neighborhood basis at $p$, and tails describe eventual membership in $\N$, giving the epsilon condition.

  \begin{lean}
    exact Metric.tendsto_atTop
  \end{lean}
\end{proof}

\begin{definition} \label{an1:def:bounded-sequence}
  A sequence is \textbf{bounded} if its range $\{u_n:n\in\N\}$ is bounded.
\end{definition}

% Lean source: Textbooks.MathematicalAnalysis1.Chapter02.ConvergentSequences.convergent_bounded
\begin{lemma} \label{an1:lem:convergent-bounded}
  If $u_n\to p$ in a metric space, then the range of $u$ is bounded.

  \begin{lean}
    theorem convergent_bounded (u : ℕ → X) (p : X) (h : Filter.Tendsto u Filter.atTop (nhds p)) :
        Bornology.IsBounded (range u)
  \end{lean}
\end{lemma}

\begin{proof}
  Use convergence with $\varepsilon=1$ to choose $N$ beyond which every term is within $1$ of $p$. Set $R=1+\sum_{n=0}^{N-1}d(u_n,p)>0$.

  \begin{lean}
    rw [tendsto_iff] at h
    apply (isBounded_iff (range u)).mpr
    intro _

    obtain ⟨N, hN⟩ := h 1 zero_lt_one

    let R := 1 + ∑ n ∈ Finset.range N, dist (u n) p

    have hsum : 0 ≤ ∑ n ∈ Finset.range N, dist (u n) p :=
      Finset.sum_nonneg (fun n _ => dist_nonneg)

    refine ⟨p, R, by
      dsimp [R]
      linarith, ?_⟩
  \end{lean}

  For $n<N$, the distance $d(u_n,p)$ is at most the finite sum defining $R-1$. For $n\geq N$, it is less than $1\leq R$. Thus every term lies in $B(p,R)$. The argument still works if the initial segment is empty.

  \begin{lean}
    rintro x ⟨n, rfl⟩
    change dist (u n) p < R
    by_cases hn : n < N

    · have hle := Finset.single_le_sum (f := fun n => dist (u n) p)
        (fun n _ => dist_nonneg) (Finset.mem_range.mpr hn)
      dsimp [R]
      linarith

    · have htail := hN n (Nat.le_of_not_gt hn)
      dsimp [R]
      linarith
  \end{lean}

\end{proof}

% Lean source: Textbooks.MathematicalAnalysis1.Chapter02.ConvergentSequences.limit_unique
\begin{lemma} \label{an1:lem:limit-unique}
  If a sequence converges to both $p$ and $q$ in a metric space, then $p=q$.

  \begin{lean}
    theorem limit_unique (u : ℕ → X) (p q : X)
        (hp : Filter.Tendsto u Filter.atTop (nhds p))
        (hq : Filter.Tendsto u Filter.atTop (nhds q)) : p = q
  \end{lean}
\end{lemma}

\begin{proof}
  If $p\ne q$, let $r=d(p,q)>0$. Beyond the larger of the two thresholds for radius $r/2$, a single term is within $r/2$ of both points. The triangle inequality would then give $d(p,q)<r$, a contradiction.

  \begin{lean}
    rw [tendsto_iff] at hp hq
    by_contra hne

    have hd : 0 < dist p q := dist_pos.mpr hne

    obtain ⟨N, hN⟩ := hp (dist p q / 2) (half_pos hd)
    obtain ⟨M, hM⟩ := hq (dist p q / 2) (half_pos hd)

    have hn := hN (max N M) (le_max_left _ _)

    have hm := hM (max N M) (le_max_right _ _)

    have ht := dist_triangle p (u (max N M)) q

    rw [dist_comm p (u (max N M))] at ht
    linarith
  \end{lean}

\end{proof}

% Lean source: Textbooks.MathematicalAnalysis1.Chapter02.ConvergentSequences.convergent_properties
\begin{proposition} \label{an1:thm:convergent-sequence-is-bounded}
  A convergent sequence is bounded and has a unique limit.

  \begin{lean}
    theorem convergent_properties (u : ℕ → X) (p : X) (hp : Filter.Tendsto u Filter.atTop (nhds p)) :
        Bornology.IsBounded (range u) ∧ ∀ q : X, Filter.Tendsto u Filter.atTop (nhds q) → p = q
  \end{lean}
\end{proposition}

\begin{proof}
  The two preceding arguments give boundedness of the range and equality of any other limit with the given one.

  \begin{lean}
    exact ⟨convergent_bounded u p hp, fun q hq => limit_unique u p q hp hq⟩
  \end{lean}

\end{proof}

% Lean source: Textbooks.MathematicalAnalysis1.Chapter02.ConvergentSequences.reciprocal_small
\begin{lemma} \label{an1:lem:reciprocal-small}
  For every $\varepsilon>0$, some $N\in\N$ satisfies $1/(n+1)<\varepsilon$ for all $n\geq N$.

  \begin{lean}
    theorem reciprocal_small (ε : ℝ) (hε : 0 < ε) :
        ∃ N : ℕ, ∀ n ≥ N, 1 / ((n : ℝ) + 1) < ε
  \end{lean}
\end{lemma}

\begin{proof}
  By the Archimedean property, choose $N>1/\varepsilon$. For $n\geq N$, positivity gives $1<N\varepsilon<(n+1)\varepsilon$. Dividing by $n+1>0$ proves the claim.

  \begin{lean}
    obtain ⟨N, hN⟩ := exists_nat_gt (1 / ε)

    refine ⟨N, ?_⟩
    intro n hn

    have hcast : (N : ℝ) ≤ n := Nat.cast_le.mpr hn

    have hprod : 1 < (N : ℝ) * ε := (div_lt_iff₀ hε).mp hN

    apply (div_lt_iff₀ (show 0 < (n : ℝ) + 1 by positivity)).mpr
    nlinarith
  \end{lean}

\end{proof}

\begin{example}
  Define $u_n=1/(n+1)$ for $n\in\N$. Then $u_n\to0$ in $\R$, but the same sequence diverges in the positive-real subspace $\R^{+}$.
  The shift makes every denominator positive under our zero-based indexing convention.

  Choose the threshold from the preceding reciprocal estimate. Since each term is positive, its distance from zero is exactly $1/(n+1)$, which is eventually below $\varepsilon$.

  If the sequence converged to a positive point $p$ in the subspace, it would converge to the same real point because the distances are unchanged. Its real limit is already zero. Uniqueness would give $p=0$, contradicting $p>0$.
\end{example}

% Lean source: Textbooks.MathematicalAnalysis1.Chapter02.ConvergentSequences.tendsto_iff_finite_exceptions
\begin{lemma} \label{an1:lem:convergesTo-iff-finite-exceptions}
  $u_n\to p$ if and only if, for every $\varepsilon>0$, the set of indices $\{n\in\N:d(u_n,p)\geq\varepsilon\}$ is finite.

  \begin{lean}
    theorem tendsto_iff_finite_exceptions (u : ℕ → X) (p : X) :
        Filter.Tendsto u Filter.atTop (nhds p) ↔ ∀ ε : ℝ, 0 < ε → {n : ℕ | ε ≤ dist (u n) p}.Finite
  \end{lean}
\end{lemma}

\begin{proof}
  If $u_n\to p$, choose its threshold $N$. Every exceptional index is less than $N$, so the exception set is contained in a finite initial segment.

  \begin{lean}
    rw [tendsto_iff]
    constructor
  \end{lean}

  Conversely, a finite set of natural indices has an upper bound $N$, including when it is empty. No index $n\geq N+1$ can be exceptional. Therefore every such term is within $\varepsilon$ of $p$, proving convergence.

  \begin{lean}
    · intro h ε hε

      obtain ⟨N, hN⟩ := h ε hε

      apply (Finset.range N).finite_toSet.subset
      intro n hn
      apply Finset.mem_range.mpr
      by_contra hnot
      exact (not_lt_of_ge hn) (hN n (Nat.le_of_not_gt hnot))

    · intro h ε hε

      obtain ⟨N, hN⟩ := (h ε hε).bddAbove

      refine ⟨N + 1, ?_⟩
      intro n hn
      by_contra hdist

      have hle := hN (show n ∈ {n : ℕ | ε ≤ dist (u n) p} from le_of_not_gt hdist)

      omega
  \end{lean}

\end{proof}

% Lean source: Textbooks.MathematicalAnalysis1.Chapter02.ConvergentSequences.sequence_at_limit_point
\begin{lemma} \label{an1:lem:sequence-at-limit-point}
  If $p\in E^{\prime}$, there is a sequence $(u_n)$ of points in $E\setminus\{p\}$ that converges to $p$.

  \begin{lean}
    theorem sequence_at_limit_point (E : Set X) (p : X) (hp : p ∈ derivedSet E) :
        ∃ u : ℕ → X, (∀ n, u n ∈ E ∧ u n ≠ p) ∧ Filter.Tendsto u Filter.atTop (nhds p)
  \end{lean}
\end{lemma}

\begin{proof}
  For each $n\in\N$, use the limit-point condition to choose $u_n\in E$, distinct from $p$, with $d(u_n,p)<1/(n+1)$.

  \begin{lean}
    rw [mem_derivedSet_iff] at hp
    have hex : ∀ n : ℕ, ∃ q ∈ E, q ≠ p ∧ dist q p < 1 / ((n : ℝ) + 1) :=
      fun n => hp _ (by positivity)

    choose u hu hne hd using hex

    refine ⟨u, fun n => ⟨hu n, hne n⟩, ?_⟩
  \end{lean}

  Given $\varepsilon>0$, the reciprocal estimate makes $1/(n+1)<\varepsilon$ for all sufficiently large $n$. The chosen distance is smaller still, so $u_n\to p$.

  \begin{lean}
    rw [tendsto_iff]
    intro ε hε

    obtain ⟨N, hN⟩ := reciprocal_small ε hε

    exact ⟨N, fun n hn => (hd n).trans (hN n hn)⟩
  \end{lean}

\end{proof}

% Lean source: Textbooks.MathematicalAnalysis1.Chapter02.ConvergentSequences.neighborhood_characterization
\begin{theorem} \label{an1:thm:neighborhood}
  A sequence converges to $p$ exactly when every positive-radius ball $B(p,\varepsilon)$ contains all but finitely many of its terms, counted by indices. Also, if $p\in E^{\prime}$, a sequence in $E\setminus\{p\}$ converges to $p$.

  \begin{lean}
    theorem neighborhood_characterization (u : ℕ → X) (p : X) :
        (Filter.Tendsto u Filter.atTop (nhds p) ↔ ∀ ε : ℝ, 0 < ε → {n : ℕ | ε ≤ dist (u n) p}.Finite) ∧
        (∀ E : Set X, p ∈ derivedSet E →
          ∃ v : ℕ → X, (∀ n, v n ∈ E ∧ v n ≠ p) ∧ Filter.Tendsto v Filter.atTop (nhds p))
  \end{lean}
\end{theorem}

\begin{proof}
  The finite-exception characterization proves the first assertion, including its converse. The construction from shrinking punctured balls proves the second assertion, with the terms explicitly required to belong to $E\setminus\{p\}$.

  \begin{lean}
    exact ⟨tendsto_iff_finite_exceptions u p, fun E hp => sequence_at_limit_point E p hp⟩
  \end{lean}

\end{proof}
